\documentclass[10pt,a4paper]{amsart}
\usepackage[margin=19mm]{geometry}
\usepackage{amsmath,amssymb,mathtools}
\usepackage[scr=rsfs,cal=euler]{mathalfa}
\usepackage{xfrac}
\usepackage[unicode,hidelinks,bookmarks=false]{hyperref}
\newcommand{\ie}{i.e.\ }
\newcommand{\cf}{cf.\ }
\newcommand{\resp}{respectively\ }
\newcommand{\Subsection}{\subsection}
\providecommand{\nobreakdash}{\nobreak\hbox{-}}
\setlength{\emergencystretch}{2em}
\multlinegap0pt
\usepackage{amssymb}
\newtheorem{lem}{Lemma}
\newtheorem{thm}{Theorem}
\title{On the Third Hankel Determinant for Inverse Coefficients of Starlike Functions: A Bernstein Polynomial Approach}
\author{Vasudevarao Allu, Shobhit Kumar}
\date{Source: arXiv:2604.23220v1 (CC BY 4.0)}
\begin{document}
\maketitle

\noindent\emph{Source and licence.} On the Third Hankel Determinant for Inverse Coefficients of Starlike Functions: A Bernstein Polynomial Approach. Version arXiv:2604.23220v1 (CC BY 4.0), distributed by arXiv under the Creative Commons Attribution 4.0 International licence, http://creativecommons.org/licenses/by/4.0/, as recorded in the arXiv licence metadata for this exact version and shown on the abstract page https://arxiv.org/abs/2604.23220v1. Full licence notice: \texttt{licenses/arxiv-2604.23220-CC-BY-4.0.txt}.\par\smallskip

\noindent\textit{Editorial scope.} Counted unit: the article's thm thm:2.1 (source0.tex lines 128--135). The article's complete original proof of it is delivered with it (source0.tex lines 137--397, one \texttt{proof} environment of 261 lines). No mathematics is written, repaired or re-derived: the statement and the proof are the article's own lines, in the article's own notation, and no retained line is substituted or re-flowed.  Retained uncounted as the source-local context the proof needs: lines 48--51; lines 103--106; lines 109--124.  Nothing was omitted from any retained span, so the package carries no omission record.  The retained lines cite 1 works, whose entries are transcribed verbatim from the article's own bibliography source in the e-print and delivered with short editorial labels recorded in the item's reference mapping.  No retained line contains a figure environment, an image-inclusion command or drawing code, so no image is delivered and the package declares no asset dependency.  Editorial formatting only, in addition: an AMS \texttt{amsart} preamble that restates the article's own declarations and macros used by the retained lines, namely the environments \texttt{lem}, \texttt{thm} on separate counters, together with the standard packages those lines need.  The retained environments print in this document as Lemma 1, Theorem 1 in order of appearance, whereas the article prints the same items with the numbering of its own full document.  The complete native source of the article as distributed in the arXiv e-print for this exact version is bundled unchanged as \texttt{sources/199148/source0.tex}.
\begin{equation}
	\label{eq:1.1}
	f(z)=z+\sum_{n=2}^{\infty}a_n z^n,\quad \text{ for } z\in\mathbb{D},
\end{equation}

\begin{equation}
	\label{eq:1.2}
	p(z)=1+\sum_{n=1}^{\infty}c_n z^n
\end{equation}

\begin{lem}{\cite{LiberaZlotkiewicz1982}}
	\label{lem:1.1}
		Let $p \in \mathcal{P}$ and be given by \eqref{eq:1.2} with $c_1 \ge 0$. Then
		\begin{align*}
			2 c_2 &= c_1^2 + \gamma \bigl(4 - c_1^2\bigr), \\[2mm]
			4 c_3 &= c_1^3 + 2\bigl(4 - c_1^2\bigr)c_1 \gamma
			- \bigl(4 - c_1^2\bigr)c_1 \gamma^2
			+ 2\bigl(4 - c_1^2\bigr)\bigl(1 - |\gamma|^2\bigr)\eta, \\[2mm]
			8 c_4 &= c_1^4
			+ \bigl(4 - c_1^2\bigr)\gamma\bigl(c_1^2(\gamma^2 - 3\gamma + 3) + 4\gamma\bigr)\\[1mm]
			&\quad- 4\bigl(4 - c_1^2\bigr)\bigl(1 - |\gamma|^2\bigr)
			\left(c_1(\gamma - 1)\eta + \overline{\gamma}\,\eta^2
			- (1 - |\eta|^2)\rho\right)
		\end{align*}
		for some $\gamma,\eta,\rho$ such that $|\gamma| \le 1$, $|\eta| \le 1$ and $|\rho| \le 1$.
	\end{lem}

	\begin{thm}
		\label{thm:2.1}
		Let $f \in \mathcal{S}^{*}$  be given by \eqref{eq:1.1}. Then
		\[
		\bigl| H(3,1)(f^{-1}) \bigr| \le 1.
		\]
		The inequality is sharp for the Koebe function.
	\end{thm}

	\begin{proof}
		Let $f \in \mathcal{S}^{*}$. Then there exists a function $p \in \mathcal{P}$ such that
		\begin{equation}\label{eq:2.1}
			z f'(z) = f(z)\,p(z).
		\end{equation}
		Note that since $f\bigl(f^{-1}(w)\bigr) = w$, it follows that
		\begin{equation}\label{eq:2.2}
			\left.
			\begin{aligned}
				A_2 &= -a_2,\\
				A_3 &= 2a_2^2 - a_3,\\
				A_4 &= 5 a_2 a_3 - 5 a_2^3 - a_4,\\
				A_5 &= 14 a_2^4 - 21 a_3 a_2^2 + 6 a_2 a_4 + 3 a_3^2 - a_5.
			\end{aligned}
			\right\}
		\end{equation}
		Comparing coefficients of \eqref{eq:2.1}, we obtain
		\begin{equation}\label{eq:2.3}
			\left.
			\begin{aligned}
				a_2 &= c_1,\\
				a_3 &= \frac{1}{2} c_2 + \frac{1}{2} c_1^2,\\
				a_4 &= \frac{1}{6} c_1^3 + \frac{1}{2} c_1 c_2 + \frac{1}{3} c_3,\\
				a_5 &= \frac{1}{24} c_1^4 + \frac{1}{4} c_2 c_1^2 + \frac{1}{3} c_1 c_3
				+ \frac{1}{8} c_2^2 + \frac{1}{4} c_4.
			\end{aligned}
			\right\}
		\end{equation}
		Using \eqref{eq:2.2} and \eqref{eq:2.3}, we obtain that
		\begin{align}\label{eq:2.4}
			144\, H(3,1)\,(f^{-1}) &=
			17 c_1^6 - 51 c_1^4 c_2 + 45 c_1^2 c_2^2 - 27 c_2^3
			+ 8 c_1^3 c_3 + 24 c_1 c_2 c_3\notag\\
			&\quad - 16 c_3^2
			- 18 c_1^2 c_4 + 18 c_2 c_4.
		\end{align}
		Since,  the class $\mathcal{S}^{*}$ is rotationally invariant. For $f \in \mathcal{S}^*$, the rotational functions 
		\[
		f_{\theta}(z):=e^{-i\theta}f(e^{i\theta}z)
		\]
		belongs to the class $\mathcal{S}^*.$ Moreover, if
		\begin{align*}
		f^{-1}(w)=w+A_{2}w^{2}+A_{3}w^{3}+A_{4}w^{4}+A_{5}w^{5}+\cdots,
		\end{align*}
		then
		\begin{align*}
		f_{\theta}^{-1}(w)=e^{-i\theta}f^{-1}(e^{i\theta}w)
		= w+e^{i\theta}A_{2}w^{2}+e^{2i\theta}A_{3}w^{3}+e^{3i\theta}A_{4}w^{4}
		+e^{4i\theta}A_{5}w^{5}+\cdots .
		\end{align*}
		Thus, we have 
		\[
		H(3,1)(f_{\theta}^{-1})=e^{6i\theta}H(3,1)(f^{-1}),
		\]
	 therefore, 
		\[
		|H(3,1)(f_{\theta}^{-1})|=|H(3,1)(f^{-1})|.
		\]
		Therefore, choosing $\theta$ suitably, we may assume that $c_{1}\in[0,2]$. Substituting the values of $c_2,c_3$ and $c_4$
		from Lemma~\ref{lem:1.1} into \eqref{eq:2.4}, and collecting
		the coefficients of $\eta$, $\eta^2$ and $\rho$, we may rewrite \eqref{eq:2.4} in the form
		\begin{equation}\label{eq:2.5}
			1152\, H(3, 1)\,(f^{-1})=\widetilde{A}+ \widetilde{B} \eta + \widetilde{C} \eta^2+ \widetilde{D}\rho
		\end{equation}
		where,
		\begin{align*}
			\widetilde{A}&=18\,c_{1}^{6}
			\;+\;
			51\,\gamma\,c_{1}^{4}\,(-4 + c_{1}^{2})
			\;+\;
			\gamma^{4}\,c_{1}^{2}\,(-4 + c_{1}^{2})^{2}
			\,+\,
			\gamma^{2}c_{1}^{2}\,(688 - 368c_{1}^{2} + 49c_{1}^{4})
			\;\\[3pt]
			&\quad+\gamma^{3}\,(-1152 + 704c_{1}^{2} - 172c_{1}^{4} + 17c_{1}^{6}), \\[3pt]
			\widetilde{B}&=4\,(-1 + |\gamma|^{2})\,c_{1}\,(-4 + c_{1}^{2})\,
			\Big( 3c_{1}^{2} + \gamma^{2}(-4 + c_{1}^{2}) + 4\gamma(5 + c_{1}^{2}) \Big),\\[3pt]
			\widetilde{C}&=-4\,(-1 + |\gamma|^{2})\,(-4 + c_{1}^{2})\,
			\Bigl(
			-8(-4 + c_{1}^{2})
			+ 8\,|\gamma|^{2}(-4 + c_{1}^{2})
			- 9\,(c_{1}^{2} + \gamma(-4 + c_{1}^{2}))\,\overline{\gamma}
			\Bigr),\\[2pt]
			\widetilde{D}&=36\,(-1 + |\gamma|^{2})(-1 + |\eta|^{2})\,(-4 + c_{1}^{2})\,
			\bigl(c_{1}^{2} + \gamma(-4 + c_{1}^{2})\bigr).
		\end{align*}
		Now, taking
		$
		x:=|\gamma|\in[0,1],\, y:=|\eta|\in[0,1].
		$
		Then, using the triangle inequality and the fact that $|\rho|\le 1$, we obtain
		\[
		1152\,\bigl|H(3,1)(f^{-1})\bigr|
		\le |\widetilde{A}|+|\widetilde{B}|\,|\eta|+|\widetilde{C}|\,|\eta|^{2}+|\widetilde{D}|\,|\rho|
		\le H(c_{1},x,y).
		\]
		For convenience, take $c_{1}=p_{1}$, where
		\begin{align}\label{eq:2.6}
			H(p_{1}, x, y) :=\;
			&\,18 p_{1}^{6}
			+ 51x\,p_{1}^{4}(4 - p_{1}^{2})
			+ x^{2} p_{1}^{2}(688 - 368 p_{1}^{2} + 49 p_{1}^{4}) \notag\\[4pt]
			& - x^{3}\!\left(-1152 + 704 p_{1}^{2} - 172 p_{1}^{4} + 17 p_{1}^{6}\right)
			+ x^{4} p_{1}^{2}(-4 + p_{1}^{2})^{2} \notag\\[4pt]
			& + 4\, y\,(1 - x^{2}) p_{1}(4 - p_{1}^{2})
			\left(3 p_{1}^{2} + x^{2}(4 - p_{1}^{2}) + 4x(5 + p_{1}^{2})\right) \notag\\[4pt]
			& + 4\,y^{2}\,(1 - x^{2})(4 - p_{1}^{2})
			\left(8(1 - x^{2})(4 - p_{1}^{2})
			+ 9(p_{1}^{2}x + x^{2}(4 - p_{1}^{2}))\right) \notag\\[4pt]
			& + 36(1 - x^{2})(1 - y^{2})(4 - p_{1}^{2})
			\left(p_{1}^{2} + x(4 - p_{1}^{2})\right).
		\end{align}
		We now determine the maximum value of \(H(p_{1},x,y)\) on the cuboid
		$
		\mathcal{D}:=\{(p_1, x, y)\in \mathbb{R}^3: 0\le p_{1}\le 2,\, 0\le x\le 1,\, 0\le y\le 1\}.
		$
		Our aim is to prove 
		\[
		\max_{\mathcal{D}} H(p_{1},x,y)=1152.
		\]
		Now, consider
		\[
		\Delta(p_{1},x,y):=1152-H(p_{1},x,y).
		\]
		A direct calculation shows that
		\[
		\Delta(p_{1},x,y)
		=(2-p_{1})(2+p_{1})(1-x)\,\Phi(p_{1},x,y),
		\]
		where
		\[
		\begin{aligned}
			\Phi(p_{1},x,y)=\;&
			-p_{1}^{4}x^{3}+16p_{1}^{4}x^{2}-33p_{1}^{4}x+18p_{1}^{4} +4p_{1}^{3}x^{3}y-12p_{1}^{3}x^{2}y-28p_{1}^{3}xy\\
			&-12p_{1}^{3}y +4p_{1}^{2}x^{3}y^{2}+4p_{1}^{2}x^{3}-68p_{1}^{2}x^{2}y^{2}-64p_{1}^{2}x^{2} -4p_{1}^{2}xy^{2}\\
			&+72p_{1}^{2}x+68p_{1}^{2}y^{2}+36p_{1}^{2}-16p_{1}x^{3}y-96p_{1}x^{2}y-80p_{1}xy \\
			&-16x^{3}y^{2}+128x^{2}y^{2}+144x^{2}+16xy^{2}+144x-128y^{2}+288 .
		\end{aligned}
		\]
		Let
		$
		p_{1}=2u,$  where $0\le u\le 1.
		$
		Then \(\Phi(2u,x,y)\) is a polynomial of degree at most \(4\) in \(u\), degree at most \(3\) in \(x\), and degree at most \(2\) in \(y\). Hence it can be written in the Bernstein form
		\[
		\Phi(2u,x,y)
		=
		\sum_{i=0}^{4}\sum_{j=0}^{3}\sum_{k=0}^{2}
		b_{ijk}\,B_{i}^{4}(u)B_{j}^{3}(x)B_{k}^{2}(y),
		\]
		where
		\[
		B_{m}^{n}(t)=\binom{n}{m}t^{m}(1-t)^{n-m},\qquad 0\le t\le 1.
		\]
		The corresponding Bernstein coefficient matrices are
		\[
		M_{0}=
		\begin{pmatrix}
			288&336&432&576\\[3mm]
			288&336&432&576\\[3mm]
			312&376&\frac{4264}{9}&608\\[3mm]
			360&456&\frac{1672}{3}&672\\[3mm]
			720&688&704&768
		\end{pmatrix},
		\]
		and the maximum entry of \(M_{0}\) is
		\[
		\max M_{0}=768.
		\]
		
		\[
		M_{1}=
		\begin{pmatrix}
			288&336&432&576\\[3mm]
			288&\frac{988}{3}&\frac{1232}{3}&528\\[3mm]
			312&\frac{1088}{3}&\frac{3880}{9}&512\\[3mm]
			348&\frac{1244}{3}&\frac{1376}{3}&480\\[3mm]
			672&576&480&384
		\end{pmatrix},
		\]
		and the maximum entry of \(M_{1}\) is
		\[
		\max M_{1}=672.
		\]
		
		\[
		M_{2}=
		\begin{pmatrix}
			160&\frac{640}{3}&\frac{1072}{3}&576\\[3mm]
			160&200&\frac{944}{3}&480\\[3mm]
			\frac{688}{3}&\frac{2440}{9}&\frac{3080}{9}&416\\[3mm]
			344&384&\frac{1112}{3}&288\\[3mm]
			768&608&352&0
		\end{pmatrix},
		\]
		and the maximum entry of \(M_{2}\) is
		\[
		\max M_{2}=768.
		\]
		Since every entry of \(M_{0}\), \(M_{1}\), and \(M_{2}\) is non-negative, and since each Bernstein basis polynomial is non-negative on \([0,1]\), it follows that
		\begin{align*}
		\Phi(2u,x,y)\ge 0
		\qquad \text{for all } (u,x,y)\in[0,1]\times[0, 1]\times[0, 1].
		\end{align*}
		Equivalently,
		\begin{align*}
		\Phi(p_{1},x,y)\ge 0
		\qquad \text{for all } (p_{1},x,y)\in[0,2]\times[0,1]\times[0,1].
		\end{align*}
		Therefore,
		\[
		\Delta(p_{1},x,y)
		=(2-p_{1})(2+p_{1})(1-x)\Phi(p_{1},x,y)\ge 0,
		\]
		because
		\[
		2-p_{1}\ge 0,\qquad 2+p_{1}>0,\qquad 1-x\ge 0
		\]
		throughout the domain. Hence
		\[
		H(p_{1},x,y)\le 1152, \qquad \text{ for } \quad 0\le p_{1}\le 2,\quad 0\le x\le 1,\quad 0\le y\le 1.
		\]
		To see that the bound $1152$ is sharp, we note that
		\[
		H(p_{1},1,y)=1152
		\qquad \text{for all } 0\le p_{1}\le 2,\ \ 0\le y\le 1,
		\]
		and also
		\[
		H(2,x,y)=1152
		\qquad \text{for all } 0\le x\le 1,\ \ 0\le y\le 1.
		\]
		Thus,
		\begin{align*}
		\max_{0\le p_{1}\le 2,\;0\le x\le 1,\;0\le y\le 1} H(p_{1},x,y)=1152.
		\end{align*}
	Therefore, we have
	\begin{align}\label{eq:2.7}
	\bigl|H(3,1)(f^{-1})\bigr|\le 1.
	\end{align}
To establish sharpness, we show that the Koebe function is extremal in this case. Consider the Koebe function
	\[
	k(z)=\frac{z}{(1-z)^{2}}=z+2z^{2}+3z^{3}+4z^{4}+\cdots .
	\]
	Its inverse has the expansion
	\[
	k^{-1}(w)=w-2w^{2}+5w^{3}-14w^{4}+42w^{5}+\cdots .
	\]
	Hence
	\[
	A_{2}=-2,\qquad A_{3}=5,\qquad A_{4}=-14,\qquad A_{5}=42,
	\]
	and therefore
	\[
	H(3,1)(k^{-1})
	=
	2A_{2}A_{3}A_{4}-A_{3}^{3}-A_{4}^{2}+A_{3}A_{5}-A_{2}^{2}A_{5}
	=1.
	\]
	Thus the equality in \eqref{eq:2.7} is attained for the Koebe function, and hence the bound attained is sharp. 
	\end{proof}

\begin{thebibliography}{99}
\bibitem[Libera]{LiberaZlotkiewicz1982} R.~J.~Libera and E.~J.~Z{\l}otkiewicz, Early coefficients of the inverse of a regular convex function, \emph{Proc. Amer. Math. Soc.} \textbf{85} (1982), 225--230.
\end{thebibliography}
\end{document}
